\documentclass[10pt]{article}

%%% Doc layout
\usepackage{fullpage} 
\usepackage{booktabs}       % professional-quality tables
\usepackage{microtype}      % microtypography
\usepackage{parskip}
\usepackage{times}

%% Hyperlinks always black, no weird boxes
\usepackage[hyphens]{url}
\usepackage[colorlinks=true,pdfborder={0 0 0}]{hyperref}

%%% Math typesetting
\usepackage{amsmath,amssymb}

%%% Write out problem statements in purple, solutions in black
\usepackage{xcolor}
\newcommand{\officialdirections}[1]{{\color{purple} #1}}

%%% Avoid automatic section numbers (we'll provide our own)
\setcounter{secnumdepth}{0}

%% --------------
%% Header
%% --------------
\usepackage{fancyhdr}
\fancyhf{}
\setlength{\headheight}{15pt}
\fancyhead[C]{CS 136 - 2026f - PS1 Submission}

\fancyfoot[C]{\thepage} % page number
\renewcommand\headrulewidth{0pt}
\pagestyle{fancy}


%% --------------
%% Begin Document
%% --------------
\begin{document}

~\\ %% add vertical space

{\Large{\bf Student Name: TODO}}

~\\ %% add vertical space
{\bf Collaboration Statement:}

Total hours spent: TODO

I consulted the following human, textbook, or AI resources:
\begin{itemize}
\item TODO
\item TODO
\item $\ldots$	
\end{itemize}

By turning this document in, I attest that I have followed the 
\href{https://www.cs.tufts.edu/cs/136/2026f/index.html#collaboration}{[CS 136 Collaboration Policy]}. All solutions represent my own work. No solution text in this document was directly provided by other humans or artificial agents. 

\tableofcontents

\newpage
\officialdirections{
\subsection*{Problem 1}
Let $\mu \in (0.0, 1.0)$ be a Beta-distributed random variable: $p(\mu) = \text{Beta}(\mu|\alpha, \beta)$.

Show that $\mathbb{E}[ \mu ] = \frac{\alpha}{\alpha + \beta}$.
}

\subsection{1: Solution}

\emph{TODO YOUR SOLUTION HERE.}


\newpage
\officialdirections{
\subsection*{2: Problem Statement}
You observe a sequence of binary random variables $X_1, \ldots, X_N$. Let $S$ be the number of positive outcomes, $S = \sum_{i=1}^N X_i$. We assume the following model, as we did in class:
\begin{align*}
p(\mu) &= \textnormal{Beta}(\mu|\alpha, \beta)\\
p(x_i|\mu) &= \textnormal{Bern}(x_i|\mu)
\end{align*}
You can assume that the hyperparameters $\alpha$ and $\beta$ have already been set. 
}

\officialdirections{\subsection*{2a: Problem Statement}
Show that the posterior predictive can be written:
\begin{equation*}
p(x_*|x_1,\ldots,x_N) = E_{p(\mu|x_1,\ldots,x_N)}[p(x_*|\mu)]
\end{equation*}
}

\subsection{2a: Solution}
TODO YOUR SOLUTION HERE



\newpage
\officialdirections{\subsection*{2b: Problem Statement}
Show that after seeing the results of the $N$ training trials, the probability of the next trial $X_*$ being positive is given by:
\begin{equation*}
p(X_* = 1|x_1,\ldots,x_N) = \frac{\alpha + s}{\alpha + \beta + N}
\end{equation*}
}

\subsection{2b: Solution}
TODO YOUR SOLUTION HERE

\newpage
\officialdirections{\subsection*{2c: Problem Statement}
Show that the probability of a positive result under the posterior predictive (i.e. your answer to 2b) can be decomposed as a weighted average of prior mean and observation mean:
\begin{equation*}
\frac{\alpha + s}{\alpha + \beta + N} = q \cdot E_{p(\mu)}[\mu] + (1-q) \cdot \frac{s}{N}
\end{equation*}
Find an expression for $q$ that always satisfies $0 \leq q \leq 1$. What happens to $q$ as $N \to \infty$, and how does that relate to consistency?
}

\subsection{2c: Solution}
TODO YOUR SOLUTION HERE



\newpage
\officialdirections{
\subsection*{3: Problem Statement}
Now let's consider a more complex model: a sequence of \textbf{words} $X_1, \ldots, X_N$ that come from a finite vocabulary of size $V$. We assume the following model, as we did in class:
\begin{align*}
p(\mu) &= \textnormal{Dir}(\mu|\alpha, \ldots, \alpha)\\
p(x_i|\mu) &= \textnormal{Cat}(x_i|\mu)
\end{align*}
\vspace{-10mm}
}
\officialdirections{\subsection*{3a: Problem Statement}
Show that after seeing the $N$ training words, the probability of the next word $X_*$ being vocabulary word $v_*$ is given by:
\begin{equation}
p(X_* = v_*|x_1,\ldots,x_N) = \frac{n_{v_*} + \alpha}{N + V\alpha}
\end{equation}
Where $n_{v_*}$ is the number of occurrences of word $v_*$ in the training dataset.
}

\subsection{3a: Solution}

TODO YOUR SOLUTION HERE

\newpage
\officialdirections{\subsection*{3b: Problem Statement}
In CA1, we'll want to compute the \textbf{marginal likelihood} or \textbf{evidence} of the observed training data which, unlike the likelihoods, integrates out the parameter $\mu$.
\begin{align*}
p(x_1,\ldots,x_N) &= \int p(\mu, x_1, \ldots, x_N) d\mu\\
&= \int p(\mu)p( x_1, \ldots, x_N|\mu) d\mu
\end{align*}


Show that the evidence has the following closed-form expression:

\begin{equation*}
p(x_1,\ldots, x_N) = \frac
	{ \Gamma(V \alpha)      \prod_{v=1}^V \Gamma( n_v + \alpha ) }
	{ \Gamma(N + V \alpha ) \prod_{v=1}^V \Gamma(\alpha)         }
\end{equation*}
}

\subsection{3b: Solution}

TODO YOUR SOLUTION HERE

\end{document}
